\providecommand{\rank}{\operatorname{r}}
\providecommand{\RR}{\mathbb R}
\providecommand{\CC}{\mathbb C}
\providecommand{\KK}{\mathbb K}

% ---------- Exercise 81 ----------
\setcounter{enumi}{80}
\item {\bfseries\boldmath
Determine a basis for each of the subspaces in Exercise 62.
}

\begin{answer}
Write $E_{ij}$ for the $3\times3$ matrix with entry $1$ in position $(i,j)$
and all other entries zero.
\begin{enumerate}
\item[$\mathbf{(a)}$] The decomposition in Exercise 62(a) gives the basis
\[
 \left(\begin{bmatrix}0&1&0\\-1&0&0\\2&0&0\end{bmatrix},
       \begin{bmatrix}0&0&1\\0&2&0\\0&0&-1\end{bmatrix}\right)
\]
of $W_1$. Independence follows from entries $(1,2)$ and $(1,3)$.
\item[$\mathbf{(c)}$] A general matrix in $W_3$ is
\[
 \begin{bmatrix}a&b&c\\5g&d&-a\\e&f&g\end{bmatrix}.
\]
Consequently a basis is
\[
 \bigl(E_{11}-E_{23},\ E_{12},\ E_{13},\ 5E_{21}+E_{33},\ E_{22},\ E_{31},\ E_{32}\bigr).
\]
The seven coefficients are uniquely determined by the entries of the matrix,
which proves independence as well as generation.
\item[$\mathbf{(d)}$] Every matrix in $W_4$ has first row $(a,a,a)$, with
its other six entries free. Therefore a basis is
\[
 \bigl(E_{11}+E_{12}+E_{13},\ E_{21},\ E_{22},\ E_{23},\ E_{31},\ E_{32},\ E_{33}\bigr).
\]
Again the seven coefficients are uniquely determined by the entries.
\end{enumerate}
The sets $W_2,W_5$ are not subspaces and have no vector space basis with the
inherited operations.
\end{answer}
